\documentclass[12pt, a4paper]{article}

%--- Paquetes ---
\usepackage[utf8]{inputenc}
\usepackage[T1]{fontenc}
\usepackage[spanish]{babel}
\usepackage{geometry}
\geometry{margin=2.5cm}
\usepackage{amsmath}
\usepackage{booktabs}
\usepackage{array}
\usepackage{xcolor}
\usepackage{listings}
\usepackage{hyperref}
\usepackage{enumitem}
\usepackage{parskip}
\usepackage{titlesec}
\usepackage{mdframed}

%--- Colores ---
\definecolor{codeblue}{RGB}{0,90,160}
\definecolor{codegray}{RGB}{245,245,245}
\definecolor{okgreen}{RGB}{0,140,0}
\definecolor{warnred}{RGB}{200,0,0}
\definecolor{warnorange}{RGB}{200,100,0}
\definecolor{juliapurple}{RGB}{100,0,170}

%--- Estilo de código Julia ---
\lstdefinelanguage{Julia}{
  keywords={using, for, in, end, if, else, function, return, true, false},
  keywordstyle=\color{juliapurple}\bfseries,
  comment=[l]{\#},
  commentstyle=\color{gray}\itshape,
  stringstyle=\color{codeblue},
  basicstyle=\ttfamily\small,
  backgroundcolor=\color{codegray},
  frame=single,
  rulecolor=\color{gray!40},
  breaklines=true,
  showstringspaces=false,
  numbers=left,
  numberstyle=\tiny\color{gray},
  numbersep=5pt,
}

%--- Comandos de estado ---
\newcommand{\cumple}{\textcolor{okgreen}{\textbf{[CUMPLE]}}}
\newcommand{\parcial}{\textcolor{warnorange}{\textbf{[CUMPLE PARCIALMENTE]}}}
\newcommand{\noimpl}{\textcolor{warnred}{\textbf{[NO IMPLEMENTADO]}}}

%--- Entorno para citas del enunciado ---
\newmdenv[
  linecolor=codeblue,
  backgroundcolor=blue!5,
  leftmargin=0pt,
  rightmargin=0pt,
  innerleftmargin=10pt,
  innerrightmargin=10pt,
  innertopmargin=6pt,
  innerbottommargin=6pt,
]{enunciadobox}

%--- Título ---
\title{
  \large Pontificia Universidad Católica de Chile \\
  \large IEE2313 --- Sistemas de Potencia \\[0.5em]
  \LARGE Revisión Punto a Punto: Enunciado vs.\ Código \texttt{ED.jl}
}
\author{Análisis del Proyecto Final}
\date{Junio 2026}

% ============================================================
\begin{document}
\maketitle
\tableofcontents
\newpage

% ============================================================
\section{Introducción}

El presente documento revisa, punto a punto, cada requisito del enunciado oficial
del Proyecto Final de IEE2313 contra el código \texttt{ED.jl} entregado como base.
Para cada sección se cita el texto del enunciado, se indica el fragmento de código
correspondiente con sus números de línea, y se justifica si el requisito se cumple,
se cumple parcialmente, o no está implementado.

El código trabaja sobre el sistema IEEE de 14 barras contenido en
\texttt{IEEE\_14\_Bus\_Proyecto.m}, con potencia base $S_\text{base} = 100$~MVA.

% ============================================================
\section{Parte I --- Caso Base (Sección 1.1)}

% ------------------------------------------------------------
\subsection{Carga del sistema y librerías}

\begin{enunciadobox}
\textit{``Su objetivo será proyectar y analizar la operación del día de mañana dentro
de este sistema, cuyos parámetros se encuentran contenidos en el archivo
\texttt{IEEE\_14\_Bus\_Proyecto.m}. Para llevarlo a cabo se deberán emplear los paquetes
\texttt{PowerSystems.jl}, \texttt{PowerFlows.jl} y \texttt{PowerSimulations.jl}
(versión 0.27.8) del ecosistema Sienna en Julia.''}
\end{enunciadobox}

\textbf{Estado:} \cumple

\begin{lstlisting}[language=Julia, firstnumber=1,
  caption={Líneas 1--20: importación de paquetes y carga del sistema}]
using PowerSystems, PowerSimulations, PowerFlows
using Dates, TimeSeries
using Ipopt
using CSV, DataFrames, Plots

file_path_static = joinpath(@__DIR__, "IEEE_14_Bus_Proyecto.m")
if !isfile(file_path_static)
    error("No se encontro el archivo .m en el directorio.")
end
sys    = System(file_path_static)
S_base = get_base_power(sys)   # S_base = 100 MVA
\end{lstlisting}

\textbf{Justificación:}
Los tres paquetes requeridos (\texttt{PowerSystems}, \texttt{PowerSimulations},
\texttt{PowerFlows}) se importan en la línea~1.
El sistema se carga desde \texttt{IEEE\_14\_Bus\_Proyecto.m} en la línea~19
mediante \texttt{System()}, y se obtiene $S_\text{base}=100$~MVA en la línea~20.
Se incluye verificación de existencia del archivo (líneas~14--16).

% ------------------------------------------------------------
\subsection{Límites Pmin/Pmax de los generadores (Tabla 1)}

\begin{enunciadobox}
\textit{``En la práctica se dispone de cinco unidades generadoras con los siguientes
límites mínimos y máximos de potencia activa [Tabla 1]:
gen-1 bus~1 $P^\text{min}=0$, $P^\text{max}=1.0$~pu;
gen-2 bus~2 $P^\text{max}=0.72$~pu;
gen-3 bus~3 $P^\text{max}=0.9$~pu;
gen-4 bus~6 $P^\text{max}=0.6$~pu;
gen-5 bus~8 $P^\text{max}=0.55$~pu.
Las potencias máximas han sido alteradas para coincidir con la Tabla~1.
Esto puede verse en el archivo \texttt{.m}.''}
\end{enunciadobox}

\textbf{Estado:} \cumple

\textbf{Justificación:}
Los límites se encuentran en la matriz \texttt{mpc.gen} del archivo
\texttt{IEEE\_14\_Bus\_Proyecto.m} y son heredados automáticamente por el código
al ejecutar \texttt{System(file\_path\_static)} en la línea~19.
La verificación contra la Tabla~1 se muestra a continuación:

\begin{table}[h!]
\centering
\begin{tabular}{ccccc}
\toprule
Enunciado & Bus & $P^\text{max}$ enunciado & $P^\text{max}$ en \texttt{.m} & ¿Coincide? \\
\midrule
gen-1 & 1 & 1.00~pu (100~MW) & 100~MW & \textcolor{okgreen}{Sí} \\
gen-2 & 2 & 0.72~pu ( 72~MW) &  72~MW & \textcolor{okgreen}{Sí} \\
gen-3 & 3 & 0.90~pu ( 90~MW) &  90~MW & \textcolor{okgreen}{Sí} \\
gen-4 & 6 & 0.60~pu ( 60~MW) &  60~MW & \textcolor{okgreen}{Sí} \\
gen-5 & 8 & 0.55~pu ( 55~MW) &  55~MW & \textcolor{okgreen}{Sí} \\
\bottomrule
\end{tabular}
\caption{Verificación de límites Pmax contra Tabla~1 del enunciado.}
\end{table}

\noindent\textbf{Nota:} El enunciado nombra los generadores gen-1 a gen-5 (por
orden de aparición). El archivo \texttt{.m} y el código los identifican como
gen-1, gen-2, gen-3, gen-6, gen-8 (por número de barra). Corresponden al mismo
conjunto de generadores.

% ------------------------------------------------------------
\subsection{Funciones de costo cuadráticas}

\begin{enunciadobox}
\textit{``Las funciones de costo de cada unidad generadora son modeladas como
funciones de costo cuadráticas y, en unidades de USD/h:}
\begin{align*}
C_1(P_1) &= 2100 + 10P_1 + 0.1P_1^2 \\
C_2(P_2) &= 7200 + 7P_2 + 0.06P_2^2 \\
C_3(P_3) &= 1500 + 6P_3 + 0.02P_3^2 \\
C_4(P_4) &= 6250 + 8P_4 + 0.07P_4^2 \\
C_5(P_5) &= 2000 + 60P_5 + 0.5P_5^2
\end{align*}
\textit{Se asumirá que todas las unidades se encuentran encendidas permanentemente.''}
\end{enunciadobox}

\textbf{Estado:} \cumple\ (para los 4 generadores que permanecen activos).

\begin{lstlisting}[language=Julia, firstnumber=65,
  caption={Líneas 65--76: asignación de funciones de costo cuadráticas}]
gens_termicos = sort!(collect(get_components(ThermalStandard, sys)),
                      by=x -> get_name(x))
# Orden resultante: gen-1, gen-2, gen-6, gen-8

# SE CAMBIARON SEGUN ENUNCIADO
costos_fijos    = [2100.0, 7200.0, 6250.0, 2000.0]          # a [USD/h]
costos_variables = [(0.1, 10.0), (0.06, 7.0), (0.07, 8.0), (0.5, 60.0)] # (c, b)

for (i, g) in enumerate(gens_termicos)
    costo_cuadratico = VariableCost(costos_variables[i])
    costo_total = ThreePartCost(costo_cuadratico, costos_fijos[i], 0.0, 0.0)
    set_operation_cost!(g, costo_total)
end
\end{lstlisting}

\textbf{Justificación:}
Tras eliminar gen-3 (línea~38), los generadores térmicos restantes ordenados
alfabéticamente son: gen-1, gen-2, gen-6, gen-8.
La función \texttt{VariableCost((c, b))} de \texttt{PowerSystems.jl} recibe
primero el coeficiente cuadrático~$c$ y luego el lineal~$b$.
La verificación coeficiente a coeficiente es:

\begin{table}[h!]
\centering
\begin{tabular}{ccccccc}
\toprule
Índice & Gen (código) & Gen (enunciado) & $a$ & $b$ & $c$ & ¿Correcto? \\
\midrule
1 & gen-1 & gen-1 & 2100 & 10.0 & 0.10 & \textcolor{okgreen}{Sí} \\
2 & gen-2 & gen-2 & 7200 &  7.0 & 0.06 & \textcolor{okgreen}{Sí} \\
3 & gen-6 & gen-4 & 6250 &  8.0 & 0.07 & \textcolor{okgreen}{Sí} \\
4 & gen-8 & gen-5 & 2000 & 60.0 & 0.50 & \textcolor{okgreen}{Sí} \\
\bottomrule
\end{tabular}
\caption{Verificación de coeficientes de costo. $C_3$ no se asigna porque
gen-3 es reemplazado por la planta solar.}
\end{table}

La suposición de unidades encendidas permanentemente se satisface mediante el
uso de \texttt{template\_economic\_dispatch()} (línea~110), que formula el ED
como un problema continuo sin variables de compromiso on/off.

% ------------------------------------------------------------
\subsection{Reemplazo de gen-3 por planta fotovoltaica}

\begin{enunciadobox}
\textit{``A partir del día de mañana, la unidad gen-3 será retirada y reemplazada
por una planta fotovoltaica de igual $P^\text{max}$ que su predecesora. Esta planta
operará con un factor de potencia unitario, sin ser capaz de regular reactivos y
con un costo variable nulo. Despachará potencia activa en función de un perfil de
irradiancia solar normalizada. La barra que la aloja no tiene la capacidad por sí
sola de regular su voltaje.''}
\end{enunciadobox}

\textbf{Estado:} \cumple\ (todos los sub-requisitos).

\begin{lstlisting}[language=Julia, firstnumber=34,
  caption={Líneas 34--59: retiro de gen-3 e incorporación de la planta solar}]
gen_termico_a_retirar      = get_component(ThermalStandard, sys, "gen-3")
max_active_power_gen_a_retirar = get_max_active_power(gen_termico_a_retirar)
bus_solar = get_bus(gen_termico_a_retirar)

remove_component!(sys, gen_termico_a_retirar)        # retira gen-3

cap_max_gen_solar = round(max_active_power_gen_a_retirar, digits=2) # 0.9 pu
set_bustype!(bus_solar, PowerSystems.ACBusTypes.PQ)  # barra pasa a PQ

gen_solar = RenewableDispatch(
    name="gen-solar",
    available=true,
    bus=bus_solar,
    active_power=0.0,
    reactive_power=0.0,
    rating=cap_max_gen_solar,                        # igual Pmax que gen-3
    prime_mover_type=PrimeMovers.PVe,
    reactive_power_limits=(min=0, max=0),            # sin regulacion de Q
    power_factor=1.0,                                # FP unitario
    operation_cost=TwoPartCost(VariableCost(0.0), 0.0), # costo nulo
    base_power=S_base)
add_component!(sys, gen_solar)
\end{lstlisting}

\textbf{Justificación sub-requisito a sub-requisito:}
\begin{enumerate}[leftmargin=*, label=\alph*)]
  \item \textbf{Igual $P^\text{max}$:} Se obtiene con \texttt{get\_max\_active\_power}
        (línea~35) y se asigna como \texttt{rating} (línea~51). Valor: 0.9~pu = 90~MW.
  \item \textbf{Factor de potencia unitario:} \texttt{power\_factor=1.0} (línea~54).
  \item \textbf{Sin regulación de reactivos:}
        \texttt{reactive\_power\_limits=(min=0,~max=0)} (línea~53). El comentario
        de esa línea está redactado de manera confusa, pero el valor implementado
        es correcto.
  \item \textbf{Costo variable nulo:}
        \texttt{TwoPartCost(VariableCost(0.0),~0.0)} (línea~55).
  \item \textbf{Barra sin control de voltaje:} La barra pasa de tipo PV a PQ
        mediante \texttt{set\_bustype!(..., ACBusTypes.PQ)} (línea~43).
  \item \textbf{Despacho según perfil de irradiancia:} Implementado en
        líneas~102--104 (ver sección~\ref{sec:perfiles}).
\end{enumerate}

% ------------------------------------------------------------
\subsection{Perfiles normalizados de demanda e irradiancia}
\label{sec:perfiles}

\begin{enunciadobox}
\textit{``se dispone de un perfil de demanda normalizado el cual debe ser considerado
para cada una de las cargas del sistema (...) la generación solar de mañana queda
determinada por la ponderación del perfil de irradiancia por la capacidad de la
unidad fotovoltaica.''}
\end{enunciadobox}

\textbf{Estado:} \cumple

\begin{lstlisting}[language=Julia, firstnumber=87,
  caption={Líneas 87--107: definición de timestamps y asignación de perfiles}]
start_time = DateTime("2024-01-01T00:00:00")
timestamps = [start_time + Hour(i) for i in 0:23]  # 24 marcas horarias

ruta_perfiles = joinpath(@__DIR__, "perfiles_normalizados.csv")
df_perfiles   = CSV.read(ruta_perfiles, DataFrame)

# Perfil de demanda -> todas las cargas
perfil_demanda_pu = df_perfiles.Demanda_normalizada
for load in get_components(PowerLoad, sys)
    ta = TimeArray(timestamps, perfil_demanda_pu)
    add_time_series!(sys, load,
        SingleTimeSeries(name="max_active_power", data=ta))
end

# Perfil de irradiancia -> generacion solar
perfil_solar_pu = df_perfiles.Irradiancia_normalizada
ta_solar = TimeArray(timestamps, perfil_solar_pu)
add_time_series!(sys, gen_solar,
    SingleTimeSeries(name="max_active_power", data=ta_solar))

transform_single_time_series!(sys, 24, Hour(1))
\end{lstlisting}

\textbf{Justificación:}
El archivo \texttt{perfiles\_normalizados.csv} contiene las columnas
\texttt{Demanda\_normalizada} e \texttt{Irradiancia\_normalizada} para las 24~horas
del día (horas~0 a~23). Se crean 24~timestamps horarios (líneas~87--88).
El perfil de demanda se aplica individualmente a \textbf{cada carga} del sistema
mediante un bucle \texttt{for} (líneas~96--99), cumpliendo el requisito de escalar
la demanda por barra según el perfil global.
El perfil solar se asigna a la planta fotovoltaica (líneas~102--104), limitando
su generación hora a hora como $P_\text{solar}(t) \leq P^\text{max}_\text{solar}
\times \text{irradiancia}(t)$.
La función \texttt{transform\_single\_time\_series!} (línea~107) convierte los
objetos \texttt{SingleTimeSeries} al formato \texttt{Deterministic} requerido
por \texttt{DecisionModel}.

% ------------------------------------------------------------
\subsection{Voltajes de consigna $\leq 1.05$~pu}

\begin{enunciadobox}
\textit{``el sistema IEEE de 14 barras ha sido modificado para que la magnitud de
los voltajes de consigna en todas las barras no supere los 1.05~pu.''}
\end{enunciadobox}

\textbf{Estado:} \cumple\ (implementado en el archivo \texttt{.m}).

\textbf{Justificación:}
Los voltajes de consigna se encuentran en la columna~$V_m$ de la matriz
\texttt{mpc.bus} del archivo \texttt{.m} y son heredados automáticamente en
la línea~19. Los valores son: 1.04, 1.03, 1.01, 1.00, 1.00, 1.025, 1.00, 1.035,
1.00, 1.00, 1.00, 1.00, 1.00, 1.00~pu para las barras~1 a~14 respectivamente.
Todos están estrictamente por debajo de 1.05~pu.

% ------------------------------------------------------------
\subsection{Escalado de demanda $+10\%$}

\begin{enunciadobox}
\textit{``los valores de potencias activas y reactivas de cada carga han sido
aumentadas en un 10\% respecto de las importadas desde el archivo \texttt{.m},
esto ya implementado en el archivo \texttt{ED.jl}.''}
\end{enunciadobox}

\textbf{Estado:} \cumple

\begin{lstlisting}[language=Julia, firstnumber=22,
  caption={Líneas 22--34: escalado de la demanda en un 10\%}]
# ################  NO MODIFICAR ESTA PONDERACION  ################
lambda_load = 1.10
cargas = collect(get_components(PowerLoad, sys))

for load in cargas
    set_active_power!(load,  get_active_power(load)  * lambda_load)
    set_reactive_power!(load, get_reactive_power(load) * lambda_load)
    
    set_max_active_power!(load, get_max_active_power(load) * lambda_load)
    set_max_reactive_power!(load, get_max_reactive_power(load) * lambda_load)
end
# ################  NO MODIFICAR ESTA PONDERACION  ################
\end{lstlisting}

\textbf{Justificación:}
Se aplica el factor $\lambda_\text{load}=1.10$ tanto a la potencia activa como
reactiva de todas las cargas del sistema, incluyendo sus límites máximos. El bloque está delimitado por comentarios
\texttt{NO MODIFICAR}, indicando que es parte del código base entregado o modificado según indicaciones.

% ------------------------------------------------------------
\subsection{Formulación del ED de 24 horas}

\begin{enunciadobox}
\textit{``para un horizonte de 24 horas se solicita determinar, en cada hora,
los puntos de despacho para cada una de las unidades generadoras a través de la
resolución de un problema de despacho económico (ED) solamente considerando
restricciones de balance de generación-demanda y potencia mínima y máxima para
los generadores.''}
\end{enunciadobox}

\textbf{Estado:} \cumple

\begin{lstlisting}[language=Julia, firstnumber=112,
  caption={Líneas 112--125: formulación, construcción y resolución del ED}]
template_ed = template_economic_dispatch()
set_network_model!(template_ed,
    NetworkModel(CopperPlatePowerModel,
                 duals=[CopperPlateBalanceConstraint],
                 use_slacks=true))

optimizer = optimizer_with_attributes(Ipopt.Optimizer, "print_level" => 3)
modelo_ed = DecisionModel(template_ed, sys,
                           optimizer=optimizer, horizon=24)

build!(modelo_ed, output_dir=ed_model_output)
solve!(modelo_ed)
\end{lstlisting}

\textbf{Justificación:}
\begin{itemize}
  \item \textbf{Horizonte 24 horas:} argumento \texttt{horizon=24} (línea~119).
  \item \textbf{Solo balance y Pmin/Pmax:} el modelo \texttt{CopperPlatePowerModel}
        impone únicamente la restricción de balance global
        $\sum P_\text{gen}(t) = \sum P_\text{dem}(t)$, sin restricciones de
        flujo en líneas. Los límites Pmin/Pmax se aplican automáticamente desde
        el sistema.
  \item \textbf{Precio marginal:} el argumento
        \texttt{duals=[CopperPlateBalanceConstraint]} habilita el cálculo del
        dual de la restricción de balance, que corresponde al costo marginal
        $\lambda$ de la energía.
\end{itemize}

% ============================================================
\section{Entregables de la Parte I}

% ------------------------------------------------------------
\subsection{Entregable 1 --- Puntos de despacho de potencia activa}

\begin{enunciadobox}
\textit{``Los puntos de despacho de potencia activa (MW) de todas las unidades
generadoras en cada hora, otorgados por los EDs.''}
\end{enunciadobox}

\textbf{Estado:} \cumple\ (CSV comentado).

\begin{lstlisting}[language=Julia, firstnumber=160,
  caption={Líneas 160--171: limpieza y presentación del despacho}]
# Limpieza de valores pequeños negativos antes de exportar
for col in names(despacho_p_print)[2:end]
    despacho_p_print[!, col] = [x < 1e-5 ? 0.0 : x for x in despacho_p_print[!, col]]
end

despacho_p_print.Demanda_Total_MW = sum.(eachrow(despacho_p_print[!, 2:end]))
CSV.write(joinpath(tablas_path, "1_despachos_generacion.csv"), despacho_p_print)
\end{lstlisting}

\textbf{Justificación:}
Se extraen las variables de despacho de generadores térmicos y solar por separado
y se unen en una tabla con columna por unidad y fila por hora. Se limpian
valores residuales del optimizador menores a $10^{-5}$ y se agrega una
columna de suma total de generación por hora. El archivo CSV se guarda en
\texttt{Tablas\_Resultados\_ED/1\_despachos\_generacion.csv} (línea~171).

% ------------------------------------------------------------
\subsection{Entregable 2 --- Costo total minimizado}

\begin{enunciadobox}
\textit{``El costo total minimizado al final del día de mañana, otorgado por los EDs.''}
\end{enunciadobox}

\textbf{Estado:} \cumple

\begin{lstlisting}[language=Julia, firstnumber=139,
  caption={Líneas 139--140: extracción del costo total}]
stats = read_optimizer_stats(res)
println(stats[1, :objective_value])
# VALOR OBTENIDO: 31147.352583930755 USD
\end{lstlisting}

\textbf{Justificación:}
Se extrae el valor de la función objetivo del optimizador mediante
\texttt{read\_optimizer\_stats}. El valor obtenido y registrado en el código
es \textbf{31\,147.35~USD}.

% ------------------------------------------------------------
\subsection{Entregable 3 --- Evolución de la demanda y satisfacción}

\begin{enunciadobox}
\textit{``La evolución de la demanda total (MW) del sistema en cada hora del día,
esclareciendo si es o no satisfecha con los despachos calculados en~1.
Identifique el o los horarios de mayor demanda.''}
\end{enunciadobox}

\textbf{Estado:} \cumple

\begin{lstlisting}[language=Julia, firstnumber=181,
  caption={Líneas 181--193: cálculo de demanda, comparación explícita y hora pico}]
demanda_base_total_mw = sum(get_active_power(load)
    for load in get_components(PowerLoad, sys)) * S_base
demanda_real_mw = demanda_base_total_mw .* df_perfiles.Demanda_normalizada
df_demanda = DataFrame(DateTime=timestamps, Demanda_Total_MW=demanda_real_mw)
CSV.write(joinpath(tablas_path, "2_demanda_total_real.csv"), df_demanda)

# Comparacion explicita: generacion total despachada vs demanda real por hora
df_comparacion = DataFrame(
    DateTime      = timestamps,
    Generacion_MW = despacho_p_print.Demanda_Total_MW,
    Demanda_MW    = demanda_real_mw,
    Diferencia_MW = despacho_p_print.Demanda_Total_MW .- demanda_real_mw
)
display(df_comparacion)
idx_max_dem = argmax(demanda_real_mw)
println("Hora de mayor demanda: $(timestamps[idx_max_dem]) -> ...(valor)... MW")
CSV.write(joinpath(tablas_path, "3_comparacion_gen_dem.csv"), df_comparacion)
\end{lstlisting}

\textbf{Justificación:}
Se calcula la demanda total en MW hora a hora (líneas~181--186) y se guarda
en \texttt{2\_demanda\_total\_real.csv}. Luego se construye un
\texttt{DataFrame} de comparación (líneas~189--193) con cuatro columnas:
\texttt{DateTime}, \texttt{Generacion\_MW} (suma del despacho de todas las
unidades), \texttt{Demanda\_MW} y \texttt{Diferencia\_MW}. Esta diferencia
verifica explícitamente si el ED satisface la demanda en cada hora (diferencia
ideal~$\approx 0$). Se identifica automáticamente la hora de mayor demanda
mediante \texttt{argmax} (línea~187). El resultado se guarda en
\texttt{3\_comparacion\_gen\_dem.csv}.

% ------------------------------------------------------------
\subsection{Entregable 4 --- Costo marginal $\lambda$ (USD/MWh)}

\begin{enunciadobox}
\textit{``La evolución del costo marginal $\lambda$ de la energía (en USD/MWh)
del sistema en cada hora del día, otorgado por los EDs.''}
\end{enunciadobox}

\textbf{Estado:} \cumple\ (CSV comentado).

\begin{lstlisting}[language=Julia, firstnumber=145,
  caption={Líneas 145--151: extracción y conversión del costo marginal}]
duals     = read_duals(res)
lambda_pu = duals["CopperPlateBalanceConstraint__System"]  # en $/pu
lambda_mw = DataFrame(DateTime  = lambda_pu.DateTime,
                      Lambda_MW = lambda_pu[!, 2] ./ S_base) # USD/MWh
display(lambda_mw)
\end{lstlisting}

\textbf{Justificación:}
El dual de la restricción \texttt{CopperPlateBalanceConstraint} es el precio
sombra de la energía, expresado en \$/pu ya que el modelo trabaja en por
unidad. La conversión $\lambda_\text{MW} = \lambda_\text{pu} / S_\text{base}$
(línea~150, con $S_\text{base}=100$~MVA) transforma correctamente las unidades
a USD/MWh. El resultado se guarda en \texttt{4\_costo\_marginal.csv} (línea~194).

% ------------------------------------------------------------
\subsection{Entregable 5 --- Perfil de voltajes AC por barra y hora}

\begin{enunciadobox}
\textit{``La evolución del perfil de magnitudes de voltajes (en pu) para cada
barra del sistema en cada hora del día, encontrado tras resolver los flujos de
potencia AC. (...) deberá para cada hora resolver un flujo de potencia AC,
considerando límites de potencia reactiva para los generadores síncronos.''}
\end{enunciadobox}

\textbf{Estado:} \cumple\ (CSV comentado).

\begin{lstlisting}[language=Julia, firstnumber=191,
  caption={Líneas 191--242: flujos de potencia AC para las 24 horas}]
sys_flujo    = deepcopy(sys)
ac_pf_solver = ACPowerFlow(check_reactive_power_limits=true)

resultados_voltaje = DataFrame(Hora=1:24)
# columna "Barra_N" por cada barra del sistema

for i in 1:24
    factor_demanda = df_perfiles.Demanda_normalizada[i]
    for l in get_components(PowerLoad, sys_flujo)
        set_active_power!(l,  base_P_load[get_name(l)] * factor_demanda)
        set_reactive_power!(l, base_Q_load[get_name(l)] * factor_demanda)
    end
    factor_solar = df_perfiles.Irradiancia_normalizada[i]
    set_active_power!(gen_solar_flujo, cap_max_gen_solar * factor_solar)
    for g in gens_termicos_flujo
        p_despacho_mw = despacho_p_print[i, get_name(g)]
        set_active_power!(g, p_despacho_mw / S_base)
    end
    res_temp = PowerFlows.solve_powerflow(ac_pf_solver, sys_flujo)
    if isa(res_temp, Dict)
        df_temp = res_temp["bus_results"]
        for b in barras
            num_bar = get_number(b)
            v_actual = df_temp[df_temp.bus_number .== num_bar, :Vm][1]
            resultados_voltaje[i, "Barra_$num_bar"] = v_actual
        end
    end
end
display(resultados_voltaje)
#CSV.write(joinpath(tablas_path, "5_perfiles_voltaje.csv"), resultados_voltaje)
\end{lstlisting}

\textbf{Justificación:}
\begin{itemize}
  \item \texttt{deepcopy(sys)} (línea~201): copia independiente del sistema
        para no alterar el modelo del ED.
  \item \texttt{check\_reactive\_power\_limits=true} (línea~202): considera
        los límites Qmin/Qmax de los generadores síncronos durante el flujo AC.
  \item El bucle de 24 iteraciones actualiza hora a hora: la demanda con el
        perfil normalizado, la generación solar con el perfil de irradiancia,
        y los generadores térmicos con sus setpoints obtenidos del ED.
  \item Se extrae la magnitud de voltaje $|V|$ de cada barra desde el campo
        \texttt{:Vm} del resultado del flujo y se almacena en
        \texttt{resultados\_voltaje}.
  \item Los resultados se guardan en \texttt{5\_perfiles\_voltaje.csv} (línea~252).
\end{itemize}

% ------------------------------------------------------------
\subsection{Análisis adicional --- Norma técnica de voltaje}

\begin{enunciadobox}
\textit{``analice si las magnitudes de voltajes de las barras de la red se
mantienen dentro de los límites establecidos por la norma técnica, es decir,
$0.95~\text{pu} \leq |V| \leq 1.05~\text{pu}$''}
\end{enunciadobox}

\textbf{Estado:} \cumple

\begin{lstlisting}[language=Julia, firstnumber=255,
  caption={Líneas 255--265: verificación automática de norma técnica}]
# Verificacion norma tecnica: 0.95 <= |V| <= 1.05 pu
println("Verificacion norma tecnica de voltajes [0.95, 1.05] pu:")
df_voltajes_largo = stack(resultados_voltaje, Not(:Hora),
                          variable_name=:Barra, value_name=:Vm)
violaciones = filter(r -> r.Vm < 0.95 || r.Vm > 1.05, df_voltajes_largo)
if nrow(violaciones) == 0
    println("Todos los voltajes se mantienen dentro del rango.")
else
    println("VIOLACIONES DE VOLTAJE ENCONTRADAS:")
    display(violaciones)
end
CSV.write(joinpath(tablas_path, "5b_violaciones_voltaje.csv"), violaciones)
\end{lstlisting}

\textbf{Justificación:}
Se convierte el \texttt{DataFrame} de voltajes de formato ancho a formato
largo con \texttt{stack} (una fila por combinación hora--barra). Luego se
filtra para obtener solo las filas donde $|V| < 0.95$ o $|V| > 1.05$~pu.
Si no hay violaciones se imprime un mensaje confirmando el cumplimiento de
la norma técnica; en caso contrario se muestran las violaciones con su hora
y barra correspondiente. El resultado se guarda en
\texttt{5b\_violaciones\_voltaje.csv}.

% ============================================================
\section{Parte II --- Caso con Contingencias (Sección 1.2)}

\begin{enunciadobox}
\textit{``Para esta parte considere exactamente lo mismo que implementó en
el Caso Base, con la diferencia de que ahora a las 21:00~hrs deberá analizar
los efectos individuales de dos posibles contingencias (criterio N$-$1):
\begin{itemize}
  \item[(a)] Salida individual de la línea 2-3 del sistema.
  \item[(b)] Caída del generador síncrono dispuesto en la barra 2.
\end{itemize}
Cualquiera de los dos eventos ocurre solamente en la hora mencionada y no
se despeja durante las horas restantes del día. Para cada contingencia
plausible se deberá mostrar el perfil de magnitudes de voltaje de cada barra
durante el día de operación. Determine cuál de las dos contingencias es la
peor justificando su respuesta.''}
\end{enunciadobox}

\textbf{Estado:} \noimpl

\begin{lstlisting}[language=Julia, firstnumber=244,
  caption={Línea 244: única referencia a las contingencias en el código}]
### PARA LAS CONTINGENCIAS HACER OTRO DEEPCOPY PARA NO ALTERAR SISTEMA ORIGINAL
\end{lstlisting}

\textbf{Justificación:}
El código solo contiene el comentario de la línea~244 indicando la estrategia
(\texttt{deepcopy}), pero no hay ninguna implementación para ninguna de las
dos contingencias. Las tareas pendientes son:
\begin{enumerate}
  \item Crear una copia independiente del sistema de flujo con
        \texttt{deepcopy}.
  \item \textbf{Contingencia (a):} desconectar la rama que une las
        barras~2 y~3, y correr flujos AC para las 24~horas con la
        contingencia activa a partir de la hora~21.
  \item \textbf{Contingencia (b):} desconectar gen-2 (la barra~2 pierde su
        generador PV), correr flujos AC desde la hora~21 con el slack
        absorbiendo la potencia faltante.
  \item Guardar y comparar los perfiles de voltaje bajo contingencia versus
        el caso base.
  \item Determinar cuál contingencia produce mayores violaciones de voltaje
        o de límites operacionales.
\end{enumerate}

% ============================================================
\section{Resumen General}

\begin{table}[h!]
\centering
\renewcommand{\arraystretch}{1.3}
\begin{tabular}{p{0.55\textwidth} c p{0.25\textwidth}}
\toprule
\textbf{Requisito del Enunciado} & \textbf{Estado} & \textbf{Líneas del código} \\
\midrule
Carga del sistema desde \texttt{.m} & \textcolor{okgreen}{Cumple} & 13--20 \\
Límites Pmin/Pmax (Tabla 1) & \textcolor{okgreen}{Cumple} & Heredado de \texttt{.m} (l.~19) \\
Unidades encendidas permanentemente & \textcolor{okgreen}{Cumple} & 110 (ED continuo) \\
$C_1$: gen-1, $a$=2100, $b$=10, $c$=0.1 & \textcolor{okgreen}{Cumple} & 68--69 (índice~1) \\
$C_2$: gen-2, $a$=7200, $b$=7, $c$=0.06 & \textcolor{okgreen}{Cumple} & 68--69 (índice~2) \\
$C_4$: gen-6, $a$=6250, $b$=8, $c$=0.07 & \textcolor{okgreen}{Cumple} & 68--69 (índice~3) \\
$C_5$: gen-8, $a$=2000, $b$=60, $c$=0.5 & \textcolor{okgreen}{Cumple} & 68--69 (índice~4) \\
Solar: igual $P^\text{max}$ que gen-3 & \textcolor{okgreen}{Cumple} & 34--35, 41, 51 \\
Solar: FP = 1, sin Q, costo = 0 & \textcolor{okgreen}{Cumple} & 53--55 \\
Solar: barra pasa a PQ & \textcolor{okgreen}{Cumple} & 43 \\
Solar: perfil de irradiancia & \textcolor{okgreen}{Cumple} & 102--104 \\
Perfil demanda a todas las cargas & \textcolor{okgreen}{Cumple} & 95--99 \\
Voltajes consigna $\leq 1.05$~pu & \textcolor{okgreen}{Cumple} & Heredado de \texttt{.m} (l.~19) \\
Demanda escalada $+10\%$ & \textcolor{okgreen}{Cumple} & 24--30 \\
ED 24 horas, solo balance & \textcolor{okgreen}{Cumple} & 107--122 \\
\textbf{E1}: Despacho P activa todas las unidades & \textcolor{okgreen}{Cumple} & 131--136, 159--160 \\
\textbf{E2}: Costo total minimizado & \textcolor{okgreen}{Cumple} & 139--140 \\
\textbf{E3}: Demanda total y comparación & \textcolor{okgreen}{Cumple} & 170--189 \\
\textbf{E4}: Costo marginal $\lambda$ USD/MWh & \textcolor{okgreen}{Cumple} & 143--148 \\
\textbf{E5}: Voltajes AC 24h por barra & \textcolor{okgreen}{Cumple} & 201--252 \\
Verificación norma técnica $[0.95, 1.05]$ & \textcolor{okgreen}{Cumple} & 254--264 \\
\textbf{Contingencia (a)}: salida línea~2-3 & \textcolor{warnred}{No implementado} & (solo l.~266) \\
\textbf{Contingencia (b)}: caída gen-2 & \textcolor{warnred}{No implementado} & (solo l.~266) \\
Comparación: cuál es la peor contingencia & \textcolor{warnred}{No implementado} & --- \\
\bottomrule
\end{tabular}
\caption{Resumen del cumplimiento de todos los requisitos del enunciado.}
\end{table}

\end{document}
